%! Linear Transformations and Matrices — Unit 8, Lesson 8.8 — Homework (Answer Key)
\documentclass[10pt]{article}
\usepackage{precalculus-article}
\usepackage{precalculus-key}   % pulls in -boxes; do NOT also load -boxes
\newcommand{\TallMath}[1]{$\displaystyle #1\rule[-1.4em]{0pt}{3.2em}$}

\begin{document}
\pageheader{Unit 8, Lesson 8.8}{Homework --- Answer Key}
\namedateperiod
\vspace{0.10in}

\begin{notesbox}
\textbf{1.} Let $A = \begin{bmatrix}4&-2\\1&3\end{bmatrix}$.
Compute $A\vec{v}$ and $A\vec{w}$ for $\vec{v}=\begin{bmatrix}1\\2\end{bmatrix}$
and $\vec{w}=\begin{bmatrix}-1\\0\end{bmatrix}$. Show your work.

\vspace{0.06in}
\ans{$A\vec{v} = \begin{bmatrix}4(1)+(-2)(2)\\1(1)+3(2)\end{bmatrix} = \begin{bmatrix}0\\7\end{bmatrix}$
\qquad
$A\vec{w} = \begin{bmatrix}4(-1)+(-2)(0)\\1(-1)+3(0)\end{bmatrix} = \begin{bmatrix}-4\\-1\end{bmatrix}$}

\vspace{0.04in}
$A\vec{v} = $ \ans{$\begin{bmatrix}0\\7\end{bmatrix}$}
\qquad
$A\vec{w} = $ \ans{$\begin{bmatrix}-4\\-1\end{bmatrix}$}
\end{notesbox}

\vspace{0.10in}

\begin{notesbox}
\textbf{2.} A linear transformation maps $\hat{\imath}\mapsto\langle2,1\rangle$
and $\hat{\jmath}\mapsto\langle-1,3\rangle$.

\vspace{0.06in}
\textbf{(a)} Write the matrix $A$ associated with this transformation.

\vspace{0.06in}
$A = $ \ans{$\begin{bmatrix}2&-1\\1&3\end{bmatrix}$}

\vspace{0.12in}
\textbf{(b)} Use $A$ to compute $L\bigl(\langle4,-2\rangle\bigr)$. Show your work.

\vspace{0.06in}
\ans{$\begin{bmatrix}2&-1\\1&3\end{bmatrix}\begin{bmatrix}4\\-2\end{bmatrix}
= \begin{bmatrix}2(4)+(-1)(-2)\\1(4)+3(-2)\end{bmatrix}
= \begin{bmatrix}10\\-2\end{bmatrix}$}

\vspace{0.04in}
$L\bigl(\langle4,-2\rangle\bigr) = $ \ans{$\begin{bmatrix}10\\-2\end{bmatrix}$}
\end{notesbox}

\vspace{0.10in}

\begin{notesbox}
\textbf{3.} Use a single matrix product to transform triangle vertices
$(0,0)$, $(2,0)$, $(1,2)$ by $A=\begin{bmatrix}0&-1\\1&0\end{bmatrix}$.

\vspace{0.06in}
$V = \begin{bmatrix}0&2&1\\0&0&2\end{bmatrix}$

\vspace{0.06in}
\ans{$AV = \begin{bmatrix}0&-1\\1&0\end{bmatrix}\begin{bmatrix}0&2&1\\0&0&2\end{bmatrix}
= \begin{bmatrix}0&0&-2\\0&2&1\end{bmatrix}$}

\vspace{0.04in}
Image vertices: \ans{$(0,0)$}, \quad \ans{$(0,2)$}, \quad \ans{$(-2,1)$}
\end{notesbox}

\vspace{0.10in}

\begin{notesbox}
\textbf{4.} Explain why the mapping
$f\bigl(\langle x,y\rangle\bigr)=\langle x+1,y+1\rangle$
is \textbf{NOT} a linear transformation. (Hint: test the zero vector.)

\vspace{0.06in}
\ans{$f(\vec{0}) = f(\langle0,0\rangle) = \langle1,1\rangle \ne \vec{0}$, so $f$ does not preserve the zero vector. A linear transformation must map $\vec{0}$ to $\vec{0}$, so $f$ cannot be linear.}
\end{notesbox}

\vspace{0.10in}

\begin{notesbox}
\textbf{5. AP-Style Multiple Choice.}
Which of the following \textbf{cannot} be a linear transformation?

\vspace{0.08in}
\begin{itemize}[leftmargin=2em, itemsep=4pt]
  \item[(A)] A transformation that maps $\vec{0}$ to $\vec{0}$
  \item[(B)] A transformation that maps every vector to twice itself
  \item[(C)] A transformation that maps $\vec{0}$ to $\langle3,0\rangle$
  \item[(D)] A transformation that rotates every vector by $45°$
\end{itemize}

\vspace{0.08in}
Answer: \ans{(C)} --- A linear transformation must always send $\vec{0}$ to $\vec{0}$.
\end{notesbox}

\end{document}
